Les coefficients d'interpolation $\delta $ a determiner sont au nombre de
quatre, deux correspondant a une interpolation sur la grille en $p$, $\delta
_{p,i+1,j+1/2}$ et $\delta _{p,i,j+1/2}$ et les deux autres sur la grille en 
$\xi ,$ $\delta _{\xi ,i+1/2,j+1}$ et $\delta _{\xi ,i+1/2,j}.$ La methode
usuelle consiste a les evaluer en imposant que sous le simple effet des
collisions, donc en l'absence de tout autre mecanisme d'acceleration (champ
electrique, ondes RF,...), la solution stationnaire donnee par le code $%
f^{\left( \infty \right) }$ soit bien une Maxwellienne $f_{M}$. Dans ce cas,
comme l'operateur de collision est de symetrie spherique, $\overline{%
\overline{D}}$ est diagonal et 
\begin{eqnarray*}
\left. S_{p}\right| _{i+1,j+1/2}^{(\infty )} &=&-D_{pp,i+1,j+1/2}\frac{%
f_{0,i+3/2,j+1/2}^{(\infty )}-f_{0,i+1/2,j+1/2}^{(\infty )}}{\Delta p_{i+1}}%
+F_{p,i+1,j+1/2}f_{0,i+1,j+1/2}^{(\infty )} \\
\left. S_{p}\right| _{i,j+1/2}^{(\infty )} &=&-D_{pp,i,j+1/2}\frac{%
f_{0,i+1/2,j+1/2}^{(\infty )}-f_{0,i-1/2,j+1/2}^{(\infty )}}{\Delta p_{i}}%
+F_{p,i,j+1/2}f_{0,i,j+1/2}^{(\infty )} \\
\left. S_{\xi }\right| _{i+1/2,j+1}^{(\infty )} &=&+D_{\xi \xi ,i+1/2,j+1}%
\sqrt{1-\xi _{j+1}^{2}}\frac{f_{0,i+1/2,j+3/2}^{(\infty %
)}-f_{0,i+1/2,j+1/2}^{(\infty )}}{p_{i+1/2}\Delta \xi _{j+1}}+F_{\xi
,i+1/2,j+1}f_{0,i+1/2,j+1}^{(\infty )} \\
\left. S_{\xi }\right| _{i+1/2,j}^{(\infty )} &=&+D_{\xi \xi ,i+1/2,j}\sqrt{%
1-\xi _{j}^{2}}\frac{f_{0,i+1/2,j+1/2}^{(\infty )}-f_{0,i+1/2,j-1/2}^{(%
\infty )}}{p_{i+1/2}\Delta \xi _{j}}+F_{\xi ,i+1/2,j}f_{0,i+1/2,j}^{(\infty %
)}
\end{eqnarray*}
soit 
\begin{eqnarray*}
\Delta p_{i+1}\left. S_{p}\right| _{i+1,j+1/2}^{(\infty )}
&=&-D_{pp,i+1,j+1/2}\left( f_{0,i+3/2,j+1/2}^{(\infty )}-f_{0,i+1/2,j+1/2}^{(%
\infty )}\right)  \\
&&+\Delta p_{i+1}F_{p,i+1,j+1/2}f_{0,i+1,j+1/2}^{(\infty )} \\
\Delta p_{i}\left. S_{p}\right| _{i,j+1/2}^{(\infty )}
&=&-D_{pp,i,j+1/2}\left( f_{0,i+1/2,j+1/2}^{(\infty )}-f_{0,i-1/2,j+1/2}^{(%
\infty )}\right)  \\
&&+\Delta p_{i}F_{p,i,j+1/2}f_{0,i,j+1/2}^{(\infty )} \\
p_{i+1/2}\Delta \xi _{j+1}\left. S_{\xi }\right| _{i+1/2,j+1}^{(\infty )}
&=&+D_{\xi \xi ,i+1/2,j+1}\sqrt{1-\xi _{j+1}^{2}}\left( f_{0,i+1/2,j+3/2}^{(%
\infty )}-f_{0,i+1/2,j+1/2}^{(\infty )}\right)  \\
&&+p_{i+1/2}\Delta \xi _{j+1}F_{\xi ,i+1/2,j+1}f_{0,i+1/2,j+1}^{(\infty )} \\
p_{i+1/2}\Delta \xi _{j}\left. S_{\xi }\right| _{i+1/2,j}^{(\infty )}
&=&+D_{\xi \xi ,i+1/2,j}\sqrt{1-\xi _{j}^{2}}\left( f_{0,i+1/2,j+1/2}^{(%
\infty )}-f_{0,i+1/2,j-1/2}^{(\infty )}\right)  \\
&&+p_{i+1/2}\Delta \xi _{j}F_{\xi ,i+1/2,j}f_{0,i+1/2,j}^{(\infty )}
\end{eqnarray*}

Cela signifie donc que 
\[
\lim_{t\rightarrow+\infty}\frac{\partial f_{0}}{\partial t}=-\nabla S=0
\]
On en deduit alors que 
\[
\left\{ 
\begin{array}{l}
\Delta
p_{i+1}F_{p,i+1,j+1/2}f_{0,i+1,j+1/2}^{(\infty)}=D_{pp,i+1,j+1/2}%
\left(f_{0,i+3/2,j+1/2}^{(\infty)}-f_{0,i+1/2,j+1/2}^{(\infty)}\right) \\ 
\Delta
p_{i}F_{p,i,j+1/2}f_{0,i,j+1/2}^{(\infty)}=D_{pp,i,j+1/2}%
\left(f_{0,i+1/2,j+1/2}^{(\infty)}-f_{0,i-1/2,j+1/2}^{(\infty)}\right) \\ 
p_{i+1/2}\Delta\xi_{j+1}F_{\xi,i+1/2,j+1}f_{0,i+1/2,j+1}^{(\infty)}=-D_{\xi%
\xi,i+1/2,j+1}\sqrt{1-\xi_{j+1}^{2}}\left(f_{0,i+1/2,j+3/2}^{(%
\infty)}-f_{0,i+1/2,j+1/2}^{(\infty)}\right) \\ 
p_{i+1/2}\Delta\xi_{j}F_{\xi,i+1/2,j}f_{0,i+1/2,j}^{(\infty)}=-D_{\xi%
\xi,i+1/2,j}\sqrt{1-\xi_{j}^{2}}\left(f_{0,i+1/2,j+1/2}^{(%
\infty)}-f_{0,i+1/2,j-1/2}^{(\infty)}\right)
\end{array}
\right.
\]
et en reportant les expressions des coefficients $\delta$ on obtient 
\[
\left\{ 
\begin{array}{l}
\Delta
p_{i+1}F_{p,i+1,j+1/2}\left(\left(1-\delta_{p,i+1,j+1/2}%
\right)f_{0,i+3/2,j+1/2}^{(\infty)}+\delta_{p,i+1,j+1/2}f_{0,i+1/2,j+1/2}^{(%
\infty)}\right) \\ 
=D_{pp,i+1,j+1/2}\left(f_{0,i+3/2,j+1/2}^{(\infty)}-f_{0,i+1/2,j+1/2}^{(%
\infty)}\right) \\ 
\Delta
p_{i}F_{p,i,j+1/2}\left(\left(1-\delta_{p,i,j+1/2}%
\right)f_{0,i+1/2,j+1/2}^{(\infty)}+\delta_{p,i,j+1/2}f_{0,i-1/2,j+1/2}^{(%
\infty)}\right) \\ 
=D_{pp,i,j+1/2}\left(f_{0,i+1/2,j+1/2}^{(\infty)}-f_{0,i-1/2,j+1/2}^{(%
\infty)}\right) \\ 
p_{i+1/2}\Delta\xi_{j+1}F_{\xi,i+1/2,j+1}\left(\left(1-\delta_{%
\xi,i+1/2,j+1}\right)f_{0,i+1/2,j+3/2}^{(\infty)}+\delta_{%
\xi,i+1/2,j+1}f_{0,i+1/2,j+1/2}^{(\infty)}\right) \\ 
=-D_{\xi\xi,i+1/2,j+1}\sqrt{1-\xi_{j+1}^{2}}\left(f_{0,i+1/2,j+3/2}^{(%
\infty)}-f_{0,i+1/2,j+1/2}^{(\infty)}\right) \\ 
p_{i+1/2}\Delta\xi_{j}F_{\xi,i+1/2,j}\left(\left(1-\delta_{\xi,i+1/2,j}%
\right)f_{0,i+1/2,j+1/2}^{(\infty)}+\delta_{\xi,i+1/2,j}f_{0,i+1/2,j-1/2}^{(%
\infty)}\right) \\ 
=-D_{\xi\xi,i+1/2,j}\sqrt{1-\xi_{j}^{2}}\left(f_{0,i+1/2,j+1/2}^{(%
\infty)}-f_{0,i+1/2,j-1/2}^{(\infty)}\right)
\end{array}
\right.
\]
soit en recombinant 
\[
\frac{f_{0,i+3/2,j+1/2}^{(\infty)}}{f_{0,i+1/2,j+1/2}^{(\infty)}}=\frac{%
D_{pp,i+1,j+1/2}+\Delta p_{i+1}F_{p,i+1,j+1/2}\delta_{p,i+1,j+1/2}}{%
D_{pp,i+1,j+1/2}-\Delta
p_{i+1}F_{p,i+1,j+1/2}\left(1-\delta_{p,i+1,j+1/2}\right)}
\]
\[
\frac{f_{0,i+1/2,j+1/2}^{(\infty)}}{f_{0,i-1/2,j+1/2}^{(\infty)}}=\frac{%
D_{pp,i,j+1/2}+\Delta p_{i}F_{p,i,j+1/2}\delta_{p,i,j+1/2}}{%
D_{pp,i,j+1/2}-\Delta p_{i}F_{p,i,j+1/2}\left(1-\delta_{p,i,j+1/2}\right)}
\]

\[
\frac{f_{0,i+1/2,j+3/2}^{(\infty)}}{f_{0,i+1/2,j+1/2}^{(\infty)}}=\frac{%
-D_{\xi\xi,i+1/2,j+1}\sqrt{1-\xi_{j+1}^{2}}+p_{i+1/2}\Delta\xi_{j+1}F_{%
\xi,i+1/2,j+1}\delta_{\xi,i+1/2,j+1}}{-D_{\xi\xi,i+1/2,j+1}\sqrt{%
1-\xi_{j+1}^{2}}-p_{i+1/2}\Delta\xi_{j+1}F_{\xi,i+1/2,j+1}\left(1-\delta_{%
\xi,i+1/2,j+1}\right)}
\]

\[
\frac{f_{0,i+1/2,j+1/2}^{(\infty)}}{f_{0,i+1/2,j-1/2}^{(\infty)}}=\frac{%
-D_{\xi\xi,i+1/2,j}\sqrt{1-\xi_{j}^{2}}+p_{i+1/2}\Delta\xi_{j}F_{%
\xi,i+1/2,j}\delta_{\xi,i+1/2,j}}{-D_{\xi\xi,i+1/2,j}\sqrt{1-\xi_{j}^{2}}%
-p_{i+1/2}\Delta\xi_{j}F_{\xi,i+1/2,j}\left(1-\delta_{\xi,i+1/2,j}\right)}
\]
que l'on peut mettre sous la forme simple 
\[
\frac{f_{0,i+3/2,j+1/2}^{(\infty)}}{f_{0,i+1/2,j+1/2}^{(\infty)}}=\frac{%
1-w_{pi+1,j+1/2}^{(u)}\delta_{p,i+1,j+1/2}}{1+w_{pi+1,j+1/2}^{(u)}\left(1-%
\delta_{p,i+1,j+1/2}\right)}
\]
\[
\frac{f_{0,i+1/2,j+1/2}^{(\infty)}}{f_{0,i-1/2,j+1/2}^{(\infty)}}=\frac{%
1-w_{pi,j+1/2}^{(u)}\delta_{p,i,j+1/2}}{1+w_{pi,j+1/2}^{(u)}\left(1-%
\delta_{p,i,j+1/2}\right)}
\]

\[
\frac{f_{0,i+1/2,j+3/2}^{(\infty)}}{f_{0,i+1/2,j+1/2}^{(\infty)}}=\frac{%
1-w_{\xi i+1/2,j+1}^{(u)}\delta_{\xi,i+1/2,j+1}}{1+w_{\xi
i+1/2,j+1}^{(u)}\left(1-\delta_{\xi,i+1/2,j+1}\right)}
\]

\[
\frac{f_{0,i+1/2,j+1/2}^{(\infty)}}{f_{0,i+1/2,j-1/2}^{(\infty)}}=\frac{%
1-w_{\xi i+1/2,j}^{(u)}\delta_{\xi,i+1/2,j}}{1+w_{\xi
i+1/2,j}^{(u)}\left(1-\delta_{\xi,i+1/2,j}\right)}
\]
avec 
\[
\left\{ 
\begin{array}{l}
w_{pi+1,j+1/2}^{(u)}=-\Delta p_{i+1}\frac{F_{p,i+1,j+1/2}}{D_{pp,i+1,j+1/2}}
\\ 
w_{pi,j+1/2}^{(u)}=-\Delta p_{i}\frac{F_{p,i,j+1/2}}{D_{pp,i,j+1/2}} \\ 
w_{\xi i+1/2,j+1}^{(u)}=+p_{i+1/2}\frac{\Delta\xi_{j+1}}{\sqrt{%
1-\xi_{j+1}^{2}}}\frac{F_{\xi,i+1/2,j+1}}{D_{\xi\xi,i+1/2,j+1}} \\ 
w_{\xi i+1/2,j}^{(u)}=+p_{i+1/2}\frac{\Delta\xi_{j}}{\sqrt{1-\xi_{j}^{2}}}%
\frac{F_{\xi,i+1/2,j}}{D_{\xi\xi,i/2,j}}
\end{array}
\right.
\]

Par ailleurs, comme $\nabla S=0,$ on en deduit naturellement que 
\[
\left\{ 
\begin{array}{l}
-D_{pp}\frac{df_{0}^{(\infty)}}{dp}+F_{p}f_{0}^{(\infty)}=0 \\ 
+\frac{\sqrt{1-\xi2}}{p}D_{\xi\xi}\frac{df_{0}^{(\infty)}}{d\xi}%
+F_{\xi}f_{0}^{(\infty)}=0
\end{array}
\right.
\]
et comme $F_{\xi}=0$ pour les collisions car aucune direction angulaire
n'est privilegiee 
\[
\left\{ 
\begin{array}{l}
\frac{\partial f_{0}^{(\infty)}}{f_{0}^{(\infty)}}=\left(\frac{F_{p}}{D_{pp}}%
\right)\partial p \\ 
\frac{\partial f_{0}^{(\infty)}}{\partial\xi}=0
\end{array}
\right.
\]
ce qui indique qu'en presence de collisions seules, $f^{(\infty)}$ ne depend
pas de $\xi.$ Donc comme $\frac{\partial}{\partial p}\rightarrow\frac{d}{dp},
$ par integration on obtient pour la premiere equation 
\[
\ln f=\int\left(\frac{F_{p}}{D_{pp}}\right)dp+Cte
\]

Les valeurs de $f_{0}^{(\infty)}$ aux points de la demi-grille $%
\left(i+3/2,j+1/2\right)$ et $\left(i+1/2,j+1/2\right)$ valent ainsi 
\[
\left\{ 
\begin{array}{l}
\ln f_{0,i+3/2,j+1/2}^{(\infty)}=\int_{0}^{p_{i+3/2}}\left(\frac{%
F_{p}\left(p,\xi_{j+1/2}\right)}{D_{pp}\left(p,\xi_{j+1/2}\right)}%
\right)dp+Cte \\ 
\ln f_{0,i+1/2,j+1/2}^{(\infty)}=\int_{0}^{p_{i+1/2}}\left(\frac{%
F_{p}\left(p,\xi_{j+1/2}\right)}{D_{pp}\left(p,\xi_{j+1/2}\right)}%
\right)dp+Cte
\end{array}
\right.
\]

On en deduit alors facilement les rapports 
\[
\left\{ 
\begin{array}{l}
\frac{f_{0,i+3/2,j+1/2}^{(\infty)}}{f_{0,i+1/2,j+1/2}^{(\infty)}}%
=\exp\left[\int_{p_{i+1/2}}^{p_{i+3/2}}\left(\frac{F_{p}\left(p,\xi_{j+1/2}%
\right)}{D_{pp}\left(p,\xi_{j+1/2}\right)}\right)dp\right] \\ 
\frac{f_{0,i+1/2,j+1/2}^{(\infty)}}{f_{0,i-1/2,j+1/2}^{(\infty)}}%
=\exp\left[\int_{p_{i-1/2}}^{p_{i13/2}}\left(\frac{F_{p}\left(p,\xi_{j+1/2}%
\right)}{D_{pp}\left(p,\xi_{j+1/2}\right)}\right)dp\right]
\end{array}
\right.
\]
et comme $f_{0}^{(\infty)}$ ne depend pas de $\xi,$ 
\[
\frac{f_{0,i+1/2,j+3/2}^{(\infty)}}{f_{0,i+1/2,j+1/2}^{(\infty)}}=\frac{%
f_{0,i+1/2,j+1/2}^{(\infty)}}{f_{0,i+1/2,j-1/2}^{(\infty)}}=1
\]

Il est facile de determiner les coefficients $\delta_{p,i+1,j+1/2}$ et $%
\delta_{p,i,j+1/2}$ en calculant l'integrale par une simple methode des
trapezes, 
\begin{eqnarray*}
\int_{p_{i+1/2}}^{p_{i+3/2}}\frac{F_{p}\left(p,\xi_{j+1/2}\right)}{%
D_{pp}\left(p,\xi_{j+1/2}\right)}dp & \approx & \frac{1}{2}\left[\frac{%
F_{p}\left(p_{i+1/2},\xi_{j+1/2}\right)}{D_{pp}\left(p_{i+1/2},\xi_{j+1/2}%
\right)}+\frac{F_{p}\left(p_{i+1},\xi_{j+1/2}\right)}{D_{pp}\left(p_{i+1},%
\xi_{j+1/2}\right)}\right]\frac{\Delta p_{i+1/2}}{2} \\
& & +\frac{1}{2}\left[\frac{F_{p}\left(p_{i+1},\xi_{j+1/2}\right)}{%
D_{pp}\left(p_{i+1},\xi_{j+1/2}\right)}+\frac{F_{p}\left(p_{i+3/2},%
\xi_{j+1/2}\right)}{D_{pp}\left(p_{i+3/2},\xi_{j+1/2}\right)}\right]\frac{%
\Delta p_{i+3/2}}{2}
\end{eqnarray*}
et 
\begin{eqnarray*}
\int_{p_{i-1/2}}^{p_{i+1/2}}\frac{F_{p}\left(p,\xi_{j+1/2}\right)}{%
D_{pp}\left(p,\xi_{j+1/2}\right)}dp & \approx & \frac{1}{2}\left[\frac{%
F_{p}\left(p_{i-1/2},\xi_{j+1/2}\right)}{D_{pp}\left(p_{i-1/2},\xi_{j+1/2}%
\right)}+\frac{F_{p}\left(p_{i},\xi_{j+1/2}\right)}{D_{pp}\left(p_{i},%
\xi_{j+1/2}\right)}\right]\frac{\Delta p_{i-1/2}}{2} \\
& & +\frac{1}{2}\left[\frac{F_{p}\left(p_{i},\xi_{j+1/2}\right)}{%
D_{pp}\left(p_{i},\xi_{j+1/2}\right)}+\frac{F_{p}\left(p_{i+1/2},\xi_{j+1/2}%
\right)}{D_{pp}\left(p_{i+1/2},\xi_{j+1/2}\right)}\right]\frac{\Delta
p_{i+1/2}}{2}
\end{eqnarray*}
puis en regroupant les termes en $p_{i+1}$ et $p_{i}$ on obtient 
\begin{eqnarray*}
\int_{p_{i+1/2}}^{p_{i+3/2}}\frac{F_{p}\left(p,\xi_{j+1/2}\right)}{%
D_{pp}\left(p,\xi_{j+1/2}\right)}dp & \approx & \frac{1}{2}\frac{%
F_{p}\left(p_{i+1},\xi_{j+1/2}\right)}{D_{pp}\left(p_{i+1},\xi_{j+1/2}\right)%
}\left(\frac{\Delta p_{i+1/2}}{2}+\frac{\Delta p_{i+3/2}}{2}\right) \\
& & +\frac{F_{p}\left(p_{i+1/2},\xi_{j+1/2}\right)}{D_{pp}\left(p_{i+1/2},%
\xi_{j+1/2}\right)}\frac{\Delta p_{i+1/2}}{4}+\frac{F_{p}\left(p_{i+3/2},%
\xi_{j+1/2}\right)}{D_{pp}\left(p_{i+3/2},\xi_{j+1/2}\right)}\frac{\Delta
p_{i+3/2}}{4} \\
& = & \frac{F_{p}\left(p_{i+1},\xi_{j+1/2}\right)}{D_{pp}\left(p_{i+1},%
\xi_{j+1/2}\right)}\frac{\Delta p_{i+1}}{2} \\
& & +\frac{F_{p}\left(p_{i+1/2},\xi_{j+1/2}\right)}{D_{pp}\left(p_{i+1/2},%
\xi_{j+1/2}\right)}\frac{\Delta p_{i+1/2}}{4}+\frac{F_{p}\left(p_{i+3/2},%
\xi_{j+1/2}\right)}{D_{pp}\left(p_{i+3/2},\xi_{j+1/2}\right)}\frac{\Delta
p_{i+3/2}}{4}
\end{eqnarray*}
soit 
\begin{eqnarray*}
\int_{p_{i+1/2}}^{p_{i+3/2}}\frac{F_{p}\left(p,\xi_{j+1/2}\right)}{%
D_{pp}\left(p,\xi_{j+1/2}\right)}dp & = & \Delta p_{i+1}\frac{%
F_{p}\left(p_{i+1},\xi_{j+1/2}\right)}{D_{pp}\left(p_{i+1},\xi_{j+1/2}\right)%
} \\
& & +\frac{F_{p}\left(p_{i+1/2},\xi_{j+1/2}\right)}{D_{pp}\left(p_{i+1/2},%
\xi_{j+1/2}\right)}\frac{\Delta p_{i+1/2}}{4} \\
& & +\frac{F_{p}\left(p_{i+3/2},\xi_{j+1/2}\right)}{D_{pp}\left(p_{i+3/2},%
\xi_{j+1/2}\right)}\frac{\Delta p_{i+3/2}}{4} \\
& & -\frac{\Delta p_{i+1}}{2}\frac{F_{p}\left(p_{i+1},\xi_{j+1/2}\right)}{%
D_{pp}\left(p_{i+1},\xi_{j+1/2}\right)}
\end{eqnarray*}

Comme 
\begin{eqnarray*}
\frac{F_{p}\left(p_{i+1},\xi_{j+1/2}\right)}{D_{pp}\left(p_{i+1},\xi_{j+1/2}%
\right)} & = & \frac{\Delta p_{i+3/2}}{\Delta p_{i+1/2}+\Delta p_{i+3/2}}%
\frac{F_{p}\left(p_{i+1/2},\xi_{j+1/2}\right)}{D_{pp}\left(p_{i+1/2},%
\xi_{j+1/2}\right)} \\
& & +\frac{\Delta p_{i+1/2}}{\Delta p_{i+1/2}+\Delta p_{i+3/2}}\frac{%
F_{p}\left(p_{i+3/2},\xi_{j+1/2}\right)}{D_{pp}\left(p_{i+3/2},\xi_{j+1/2}%
\right)}
\end{eqnarray*}
soit 
\[
\frac{F_{p}\left(p_{i+1},\xi_{j+1/2}\right)}{D_{pp}\left(p_{i+1},\xi_{j+1/2}%
\right)}=\frac{1}{2\Delta p_{i+1}}\left[\Delta p_{i+3/2}\frac{%
F_{p}\left(p_{i+1/2},\xi_{j+1/2}\right)}{D_{pp}\left(p_{i+1/2},\xi_{j+1/2}%
\right)}+\Delta p_{i+1/2}\frac{F_{p}\left(p_{i+3/2},\xi_{j+1/2}\right)}{%
D_{pp}\left(p_{i+3/2},\xi_{j+1/2}\right)}\right]
\]
on trouve apres avoir regroupe les termes 
\begin{eqnarray*}
\int_{p_{i+1/2}}^{p_{i+3/2}}\frac{F_{p}\left(p,\xi_{j+1/2}\right)}{%
D_{pp}\left(p,\xi_{j+1/2}\right)}dp & = & \Delta p_{i+1}\frac{%
F_{p}\left(p_{i+1},\xi_{j+1/2}\right)}{D_{pp}\left(p_{i+1},\xi_{j+1/2}\right)%
} \\
& & +\frac{\Delta p_{i+3/2}-\Delta p_{i+1/2}}{4}\left[\frac{%
F_{p}\left(p_{i+3/2},\xi_{j+1/2}\right)}{D_{pp}\left(p_{i+3/2},\xi_{j+1/2}%
\right)}-\frac{F_{p}\left(p_{i+1/2},\xi_{j+1/2}\right)}{D_{pp}%
\left(p_{i+1/2},\xi_{j+1/2}\right)}\right]
\end{eqnarray*}

En jouant sur les indices on a de meme pour la seconde integrale

\begin{eqnarray*}
\int_{p_{i-1/2}}^{p_{i+1/2}}\frac{F_{p}\left(p,\xi_{j+1/2}\right)}{%
D_{pp}\left(p,\xi_{j+1/2}\right)}dp & \approx & =\Delta p_{i}\frac{%
F_{p}\left(p_{i},\xi_{j+1/2}\right)}{D_{pp}\left(p_{i},\xi_{j+1/2}\right)} \\
& & +\frac{\Delta p_{i+1/2}-\Delta p_{i-1/2}}{4}\left[\frac{%
F_{p}\left(p_{i+1/2},\xi_{j+1/2}\right)}{D_{pp}\left(p_{i+1/2},\xi_{j+1/2}%
\right)}-\frac{F_{p}\left(p_{i-1/2},\xi_{j+1/2}\right)}{D_{pp}%
\left(p_{i-1/2},\xi_{j+1/2}\right)}\right]
\end{eqnarray*}

En egalant les deux expressions des rapports $\frac{f_{0,i+3/2,j+1/2}^{(%
\infty)}}{f_{0,i+1/2,j+1/2}^{(\infty)}}$ on obtient facilement 
\[
\delta_{p,i+1,j+1/2}=\frac{1}{w_{pi+1,j+1/2}^{(u)}}-\frac{1}{%
\exp\left[w_{pi+1,j+1/2}^{(u)}+w_{pi+1,j+1/2}^{(nu)}\right]-1}
\]
avec 
\[
w_{pi+1,j+1/2}^{(nu)}=-\frac{\Delta p_{i+3/2}-\Delta p_{i+1/2}}{4}\left[%
\frac{F_{p}\left(p_{i+3/2},\xi_{j+1/2}\right)}{D_{pp}\left(p_{i+3/2},%
\xi_{j+1/2}\right)}-\frac{F_{p}\left(p_{i+1/2},\xi_{j+1/2}\right)}{%
D_{pp}\left(p_{i+1/2},\xi_{j+1/2}\right)}\right]
\]

De meme 
\[
\delta_{p,i,j+1/2}=\frac{1}{w_{pi,j+1/2}^{(u)}}-\frac{1}{\exp%
\left[w_{pi,j+1/2}^{(u)}+w_{pi,j+1/2}^{(nu)}\right]-1}
\]
avec 
\[
w_{pi,j+1/2}^{(nu)}=-\frac{\Delta p_{i+1/2}-\Delta p_{i-1/2}}{4}\left[\frac{%
F_{p}\left(p_{i+1/2},\xi_{j+1/2}\right)}{D_{pp}\left(p_{i+1/2},\xi_{j+1/2}%
\right)}-\frac{F_{p}\left(p_{i-1/2},\xi_{j+1/2}\right)}{D_{pp}%
\left(p_{i-1/2},\xi_{j+1/2}\right)}\right]
\]

Dans le cas d'une grille uniforme, $w_{pi,j+1/2}^{(nu)}=0,$ et on retrouve
bien la forme usuelle de la fonction de Chang et Cooper 
\[
g\left(x\right)=\frac{1}{x}-\frac{1}{e^{x}-1}
\]
avec $x=w_{p}^{(u)}.$ Un developpement limite donne $g\left(x%
\right)=0.5-x/12+x^{3}/720+{\rm O}\left(x^{4}\right)$ lorsque $x\ll1$ avec
la limite usuelle $\lim_{x\rightarrow0}g\left(x\right)=\frac{1}{2}.$ A noter
qu'il faut introduire une fonction generalisee de Chang et Cooper pour le
cas de la grille non uniforme, sous la forme simple de deux variables 
\[
g\left(x,y\right)=\frac{1}{x}-\frac{1}{e^{x+y}-1}
\]
le terme $y=w_{p}^{(nu)}$ etant nul lorsque la grille est uniforme.

Avec cette definition, 
\[
\frac{f_{0,i+3/2,j+1/2}^{(\infty)}}{f_{0,i+1/2,j+1/2}^{(\infty)}}=\frac{%
1-w_{pi+1,j+1/2}^{(u)}\left[\frac{1}{w_{pi+1,j+1/2}^{(u)}}-\frac{1}{%
\exp\left[w_{pi+1,j+1/2}^{(u)}+w_{pi+1,j+1/2}^{(nu)}\right]-1}\right]}{%
1+w_{pi+1,j+1/2}^{(u)}\left(1-\left[\frac{1}{w_{pi+1,j+1/2}^{(u)}}-\frac{1}{%
\exp\left[w_{pi+1,j+1/2}^{(u)}+w_{pi+1,j+1/2}^{(nu)}\right]-1}\right]\right)}
\]
soit 
\[
\frac{f_{0,i+3/2,j+1/2}^{(\infty)}}{f_{0,i+1/2,j+1/2}^{(\infty)}}=\frac{%
\frac{w_{pi+1,j+1/2}^{(u)}}{\exp%
\left[w_{pi+1,j+1/2}^{(u)}+w_{pi+1,j+1/2}^{(nu)}\right]-1}}{%
w_{pi+1,j+1/2}^{(u)}+\frac{w_{pi+1,j+1/2}^{(u)}}{\exp%
\left[w_{pi+1,j+1/2}^{(u)}+w_{pi+1,j+1/2}^{(nu)}\right]-1}}
\]
qui se simplifie et devient 
\[
\frac{f_{0,i+3/2,j+1/2}^{(\infty)}}{f_{0,i+1/2,j+1/2}^{(\infty)}}%
=\exp\left[-w_{pi+1,j+1/2}^{(u)}-w_{pi+1,j+1/2}^{(nu)}\right]
\]

Comme pour les collisions, on a la relation simple 
\[
\frac{F_{p}\left(p,\xi\right)}{D_{pp}\left(p,\xi\right)}\cong-p
\]
puisque la distribution $f^{(\infty)}$ est une Maxwellienne $f_{M},$ on
trouve donc

\[
\frac{f_{0,i+3/2,j+1/2}^{(\infty)}}{f_{0,i+1/2,j+1/2}^{(\infty)}}%
=\exp\left[-\Delta p_{i+1}p_{i+1}-\frac{\Delta p_{i+3/2}-\Delta p_{i+1/2}}{4}%
\left(p_{i+3/2}-p_{i+1/2}\right)\right]
\]
En utilisant les egalites 
\[
\left\{ 
\begin{array}{l}
p_{i+3/2}=p_{i+1}+\frac{\Delta p_{i+3/2}}{2} \\ 
p_{i+1/2}=p_{i+1}-\frac{\Delta p_{i+1/2}}{2}
\end{array}
\right.
\]
il est facile de montrer que 
\[
p_{i+1}=\frac{1}{2}\left(p_{i+3/2}+p_{i+1/2}\right)-\frac{1}{2}\left(\frac{%
\Delta p_{i+3/2}}{2}-\frac{\Delta p_{i+1/2}}{2}\right)
\]
et 
\[
\Delta p_{i+1}=\frac{1}{2}\left(\Delta p_{i+3/2}+\Delta
p_{i+1/2}\right)=p_{i+3/2}-p_{i+1/2}
\]
et en reportant ces relations dans l'expression precedente on trouve 
\begin{eqnarray*}
& & -\Delta p_{i+1}p_{i+1}-\frac{\Delta p_{i+3/2}-\Delta p_{i+1/2}}{4}%
\left(p_{i+3/2}-p_{i+1/2}\right) \\
& = & -\left(p_{i+3/2}-p_{i+1/2}\right)\left[\frac{1}{2}%
\left(p_{i+3/2}+p_{i+1/2}\right)-\frac{1}{4}\left(\Delta p_{i+3/2}-\Delta
p_{i+1/2}\right)\right] \\
& & -\frac{\Delta p_{i+3/2}-\Delta p_{i+1/2}}{4}\left(p_{i+3/2}-p_{i+1/2}%
\right) \\
& = & -\frac{1}{2}\left(p_{i+3/2}+p_{i+1/2}\right)\left(p_{i+3/2}-p_{i+1/2}%
\right) \\
& = & -\frac{1}{2}\left(p_{i+3/2}^{2}-p_{i+1/2}^{2}\right)
\end{eqnarray*}
soit 
\[
\frac{f_{0,i+3/2,j+1/2}^{(\infty)}}{f_{0,i+1/2,j+1/2}^{(\infty)}}=\exp\left[-%
\frac{p_{i+3/2}^{2}-p_{i+1/2}^{2}}{2}\right]
\]
ce qui correspond bien a une Maxwellienne numerique. Les erreurs ne se
propagent donc pas avec ce type de ponderation.

Pour les termes angulaires, le calcul est tres simple et des relations
precedentes on en deduit directement que $w_{\xi i+1/2,j+1}=w_{\xi i+1/2,j}=0
$ car $F_{\xi}=0$. Ainsi $\delta_{\xi,i+1/2,j+1}$ peut etre defini de
maniere arbitraire. La fa\c{c} on la plus naturelle est d'effectuer une
interpolation lineaire entre $f_{0,i+1/2,j+3/2}^{(\infty)}$ et $%
f_{0,i+1/2,j+3/2}^{(\infty)}$ ainsi qu'entre $f_{0,i+1/2,j+1/2}^{(\infty)}$
et $f_{0,i+1/2,j-1/2}^{(\infty)}$ selon les relations 
\[
\left\{ 
\begin{array}{l}
f_{0,i+1/2,j+1}^{(\infty)}=\left(\frac{\Delta\xi_{j+1/2}}{%
\Delta\xi_{j+1/2}+\Delta\xi_{j+3/2}}\right)f_{0,i+1/2,j+3/2}^{(\infty)}+%
\left(\frac{\Delta\xi_{j+3/2}}{\Delta\xi_{j+1/2}+\Delta\xi_{j+3/2}}%
\right)f_{0,i+1/2,j+1/2}^{(\infty)} \\ 
f_{0,i+1/2,j}^{(\infty)}=\left(\frac{\Delta\xi_{j-1/2}}{\Delta\xi_{j-1/2}+%
\Delta\xi_{j+1/2}}\right)f_{0,i+1/2,j+1/2}^{(\infty)}+\left(\frac{%
\Delta\xi_{j+1/2}}{\Delta\xi_{j-1/2}+\Delta\xi_{j+1/2}}%
\right)f_{0,i+1/2,j-1/2}^{(\infty)}
\end{array}
\right.
\]
en tenant compte des pas de grille, ce qui donne 
\[
\left\{ 
\begin{array}{l}
\delta_{\xi,i+1/2,j+1}=\frac{\Delta\xi_{j+3/2}}{\Delta\xi_{j+1/2}+\Delta%
\xi_{j+3/2}} \\ 
\delta_{\xi,i+1/2,j}=\frac{\Delta\xi_{j+1/2}}{\Delta\xi_{j-1/2}+\Delta%
\xi_{j+1/2}}
\end{array}
\right.
\]
On retrouve bien que pour une grille uniforme, $\delta_{\xi,i+1/2,j+1}=%
\delta_{\xi,i+1/2,j}=1/2.$

En resume, pour les coefficients d'interpolation $\delta,$ on obtient pour
la grille en $p$%
\[
\delta_{p}\rightarrow\left\{ 
\begin{array}{l}
\delta_{p,i+1,j+1/2}=\frac{1}{w_{pi+1,j+1/2}^{(u)}}-\frac{1}{%
\exp\left[w_{pi+1,j+1/2}^{(u)}+w_{pi+1,j+1/2}^{(nu)}\right]-1} \\ 
\delta_{p,i,j+1/2}=\frac{1}{w_{pi,j+1/2}^{(u)}}-\frac{1}{\exp%
\left[w_{pi,j+1/2}^{(u)}+w_{pi,j+1/2}^{(nu)}\right]-1}
\end{array}
\right.
\]
avec 
\[
w_{p}^{(u)}\rightarrow\left\{ 
\begin{array}{l}
w_{pi+1,j+1/2}^{(u)}=-\Delta p_{i+1}\frac{F_{p,i+1,j+1/2}}{D_{pp,i+1,j+1/2}}
\\ 
w_{pi,j+1/2}^{(u)}=-\Delta p_{i}\frac{F_{p,i,j+1/2}}{D_{pp,i,j+1/2}}
\end{array}
\right.
\]
et 
\[
w_{p}^{(nu)}\rightarrow\left\{ 
\begin{array}{l}
w_{pi+1,j+1/2}^{(nu)}=-\frac{\Delta p_{i+3/2}-\Delta p_{i+1/2}}{4}\left[%
\frac{F_{p}\left(p_{i+3/2},\xi_{j+1/2}\right)}{D_{pp}\left(p_{i+3/2},%
\xi_{j+1/2}\right)}-\frac{F_{p}\left(p_{i+1/2},\xi_{j+1/2}\right)}{%
D_{pp}\left(p_{i+1/2},\xi_{j+1/2}\right)}\right] \\ 
w_{pi,j+1/2}^{(nu)}=-\frac{\Delta p_{i+1/2}-\Delta p_{i-1/2}}{4}\left[\frac{%
F_{p}\left(p_{i+1/2},\xi_{j+1/2}\right)}{D_{pp}\left(p_{i+1/2},\xi_{j+1/2}%
\right)}-\frac{F_{p}\left(p_{i-1/2},\xi_{j+1/2}\right)}{D_{pp}%
\left(p_{i-1/2},\xi_{j+1/2}\right)}\right]
\end{array}
\right.
\]
ainsi que pour la grille en $\xi$%
\[
\delta_{\xi}\rightarrow\left\{ 
\begin{array}{l}
\delta_{\xi,i+1/2,j+1}=\frac{\Delta\xi_{j+3/2}}{\Delta\xi_{j+1/2}+\Delta%
\xi_{j+3/2}} \\ 
\delta_{\xi,i+1/2,j}=\frac{\Delta\xi_{j+1/2}}{\Delta\xi_{j-1/2}+\Delta%
\xi_{j+1/2}}
\end{array}
\right.
\]

A noter que comme on ne considere que le cas des collisions pour definir les
coefficients $\delta,$ les dynamiques en $p$ et en $\xi$ sont decouplees. En
consequence, les rapports $\frac{F_{p}}{D_{pp}}$ ne dependent que de $p$, et
c'est pour cela qu'il est possible de prendre les memes valeurs en
differents points de grille en $\xi.$

Les calculs effectues dans le cas ou seules les collisions interviennent
sont usuellement etendus aux regimes en presence de mecanismes
d'acceleration (champ electrique, ondes RF) en generalisant les expressions
des coefficients de ponderations $\delta$ selon la regle simple 
\[
\left.\frac{F_{p}}{D_{pp}}\right|_{effectif}=\frac{\sum_{m}F_{p}^{(m)}}{%
\sum_{m}D_{pp}^{(m)}}
\]
ou $\sum_{m}$ est la somme sur tous les mecanismes $m$ en jeu. Ceci permet
de garder la meme syntaxe et ainsi 
\[
w_{pi+1,j+1/2}^{(nu)}=-\frac{\Delta p_{i+3/2}-\Delta p_{i+1/2}}{4}%
\left[\left.\frac{F_{p}}{D_{pp}}\right|_{effectif}^{i+3/2,j+1/2}-\left.\frac{%
F_{p}}{D_{pp}}\right|_{effectif}^{i+1/2,j+1/2}\right]
\]
et 
\[
w_{pi,j+1/2}^{(nu)}=-\frac{\Delta p_{i+1/2}-\Delta p_{i-1/2}}{4}\left[\left.%
\frac{F_{p}}{D_{pp}}\right|_{effectif}^{i+1/2,j+1/2}-\left.\frac{F_{p}}{%
D_{pp}}\right|_{effectif}^{i-1/2,j+1/2}\right]
\]

Dans ce cas evidemment, il apparait une dependence en fonction de $\xi,$
mais celle-ci est sans importance car dans la zone ou les flux sont modifies
la Maxwellienne n'est plus la solution locale. La valeur de la fonction de
Chang et Cooper tend ainsi vers la valeur $\delta_{p}\simeq\frac{1}{2}.$ On
remplace tous les coefficients des tableaux ci-dessus par leurs valeurs
effectives, comme indique.

A noter que les expressions de $w_{p}^{(nu)}$ et $\delta_{\xi,}$ font
apparaitre des termes aux points de grille en $i-1/2,$ $i+3/2$ et en $j-1/2,$
$j+3/2.$ La fa\c{c}on usuelle de les prendre en compte dans les calculs
consiste a effectuer un decalage d'indice, sous la forme 
\[
\left(i+1/2,j+1/2\right)\in\left\{
\left[1:n_{p}\right],\left[1:n_{\xi}\right]\right\} 
\]
et 
\[
\left\{ 
\begin{array}{l}
\left(i+3/2,j+1/2\right)\in\left\{
\left[2:n_{p},X\right],\left[1:n_{\xi}\right]\right\} \\ 
\left(i-1/2,j+1/2\right)\in\left\{
\left[X,1:n_{p}-1\right],\left[1:n_{\xi}\right]\right\} \\ 
\left(i+1/2,j+3/2\right)\in\left\{
\left[1:n_{p}\right],\left[2:n_{\xi},X\right]\right\} \\ 
\left(i+1/2,j-1/2\right)\in\left\{
\left[1:n_{p}\right],\left[X,1:n_{\xi}-1\right]\right\}
\end{array}
\right.
\]
Ainsi, toutes les quantites ont la meme dimension $\left(n_{p},n_{\xi}%
\right),$ ce qui est indispensable en utilisant MatLab pour une description
matricielle. Par ailleurs, les conditions aux limites aux extremites du
domaine sont naturellement satisfaites. On peut donc attribuer a la variable
une valeur $X$une valeur quelconque.
